\documentclass[11pt,a4paper]{article}

\usepackage[T1]{fontenc}
\usepackage[utf8]{inputenc}
\usepackage{lmodern}
\usepackage[margin=2.5cm]{geometry}
\usepackage{amsmath}
\usepackage{graphicx}
\usepackage[hidelinks]{hyperref}

\title{A Clear Title for Your Article}
\author{First Author \and Second Author}
\date{\today}

\begin{document}

\maketitle

\begin{abstract}
State the problem, the approach and the main result in a few sentences.
Keep it self-contained: a reader should understand the contribution
without opening the paper.
\end{abstract}

\section{Introduction}
Introduce the problem and why it matters. Cite prior work with
\verb|\cite|, for example~\cite{knuth1984texbook}.

\section{Method}
Describe what you did. Equations are numbered automatically:
\begin{equation}
  \label{eq:euler}
  e^{i\pi} + 1 = 0.
\end{equation}
Refer to them with \verb|\eqref|, as in~\eqref{eq:euler}.

\section{Results}
Report what you found. Place figures in this folder and include them with
\verb|\includegraphics|.

\section{Conclusion}
Summarise the contribution and name the open questions~\cite{lamport1994latex}.

\bibliographystyle{plain}
\bibliography{refs}

\end{document}
